\chapter{当前没实现什么，以及为什么别误读研究报告}

\section{本章要解决的困惑}

历史报告里有很多内容：Delta、Energy、OFI、MLOFI、动态 L2 path、深度学习、maker、容量曲线、Pareto。它们不等于当前 runtime。

\section{当前 runtime 已经使用}

\begin{lstlisting}
TFI
trade_window_count
past_event_25_bps
frames_since_mid_change
spread / entry_cross_bps / prior-day q70
R5
cell = 00 / 10 / 01 / 11
membership_set / base_weight
gamma / overlay
3x capacity ledger
fixed60 / simple exit
\end{lstlisting}

\section{当前 runtime 没有使用}

\begin{lstlisting}
Delta_bps / Energy_bps control
OFI / MLOFI entry trigger
dynamic L2 path prediction
deep learning
maker strategy
L2 capacity curve
path distribution model
\end{lstlisting}

\warningline{\code{Delta_bps} / \code{Energy_bps} 没进入 CC runtime 控制。\code{OFI} / \code{MLOFI} 不是当前 entry trigger。}

\section{为什么这不是否定研究}

研究报告有价值，因为它们能提出下一步问题。但工程教材必须区分：

\begin{itemize}[leftmargin=2em]
  \item 已经在 Bot/fast runtime 中执行的规则；
  \item 只在 diagnostic 中出现的 label；
  \item 还没进策略的候选因子；
  \item 未来可能做的工程扩展。
\end{itemize}

否则新人会以为系统已经很复杂，结果连当前状态机都看不懂。

\section{当前最大缺口}

当前最大缺口不是再把文档写得更玄，而是建立更严肃的预测对象：

\begin{lstlisting}
entry state X_t
-> maximum favorable move
-> time to maximum
-> time back to 0bps
-> executable return at multiple horizons
-> path class distribution
\end{lstlisting}

这些对象建立前，不要急着说深度学习，也不要先谈容量。

\checkpoint{本书最终结论很朴素：当前系统是一个简单策略加严肃工程边界。先掌握它，再研究更复杂的东西。}
